% Slide 1: Tiêu đề
% Slide 2: Mục lục
\begin{frame}{Nội dung bài thuyết trình}
    \tableofcontents
\end{frame}

\section{Công cụ: Stack \& Queue}
% Slide 3: Stack
\begin{frame}[fragile]{DFS}{1. Ngăn xếp (Stack) - Cơ chế cho DFS}
    \begin{columns}
        \begin{column}{0.55\textwidth}
            \begin{block}{Nguyên lý LIFO (Last In, First Out)}
                Phần tử vào \textbf{sau cùng} sẽ được lấy ra \textbf{đầu tiên}.
            \end{block}
            \begin{itemize}
                \item Tương tự như một chồng đĩa.
                \item Dùng để ghi nhớ các đỉnh cần quay lui trong DFS.
                \item Thao tác chính: Push (Thêm) và Pop (Lấy ra).
            \end{itemize}
        \end{column}
        \begin{column}{0.4\textwidth}
            \centering \textbf{Mô phỏng Stack}\\
            \begin{tikzpicture}
                \draw[thick] (0,3) -- (0,0) -- (1.5,0) -- (1.5,3);
                \node<1->[stack_box] at (0.75, 0.35) {Đỉnh 3};
                \node<2->[stack_box] at (0.75, 1.05) {Đỉnh 1};
                \node<3->[stack_box, fill=orange!30] at (0.75, 1.75) {Lấy ra 1};
                \draw<3->[-Stealth, red, thick] (1.6, 1.75) -- (2.5, 1.75) node[right]{Pop};
            \end{tikzpicture}
        \end{column}
    \end{columns}
\end{frame}

% Slide 4: Queue
\begin{frame}{2. Hàng đợi (Queue) - Cơ chế cho BFS}
    \begin{columns}
        \begin{column}{0.6\textwidth}
            \begin{block}{Nguyên lý FIFO (First In, First Out)}
                Phần tử vào \textbf{trước} sẽ được phục vụ \textbf{trước}.
            \end{block}
            \begin{itemize}
                \item Tương tự như xếp hàng mua vé.
                \item Giúp BFS duyệt theo từng lớp khoảng cách.
                \item Thao tác chính: Enqueue (Vào hàng) và Dequeue (Ra hàng).
            \end{itemize}
        \end{column}
        \begin{column}{0.35\textwidth}
            \centering \textbf{Mô phỏng Queue}\\
            \begin{tikzpicture}
                \draw[thick] (0,1) -- (3,1) (0,0) -- (3,0);
                \node<1->[rectangle, draw, fill=blue!10, minimum width=0.8cm] at (0.5, 0.5) {3};
                \node<1->[rectangle, draw, fill=blue!10, minimum width=0.8cm] at (1.5, 0.5) {1};
                \node<2->[rectangle, draw, fill=orange!30, minimum width=0.8cm] at (2.5, 0.5) {0};
                \draw<2->[-Stealth, thick] (3.1, 0.5) -- (3.8, 0.5) node[right]{Out};
                \draw[Stealth-, thick] (-0.8, 0.5) -- (-0.1, 0.5) node[left]{In};
            \end{tikzpicture}
        \end{column}
    \end{columns}
\end{frame}

\section{Thuật toán DFS}
% Slide 5: DFS Theory
\begin{frame}{3. DFS: Tìm kiếm theo chiều sâu}
    \begin{itemize}
        \item \textbf{Tư tưởng:} Bắt đầu từ một đỉnh nguồn, đi sâu nhất có thể theo các cạnh cho đến khi gặp ngõ cụt hoặc đỉnh đã duyệt thì mới quay lui.
        \item \textbf{Đặc điểm:} Duyệt hết các nhánh trước khi chuyển sang nhánh lân cận khác.
        \item \textbf{Độ phức tạp:} Thời gian $O(V + E)$, Không gian $O(V)$.
    \end{itemize}
\end{frame}

% Slide 6: DFS Animation (Đã nâng cấp hiệu ứng)
\begin{frame}{Mô phỏng từng bước DFS}
    \begin{columns}
        \begin{column}{0.5\textwidth}
            \centering
            \begin{tikzpicture}[scale=1.2, every node/.style={node_base}]
                % Định nghĩa các đỉnh với hiệu ứng mờ/hiện dần
                \node<1>[node_current] (0) at (0,0) {0};
                \node<2->[node_visited] (0) at (0,0) {0};

                \node<1>[node_base, opacity=0.3] (1) at (1.5, 1) {1}; % Mờ khi chưa tới
                \node<2>[node_current] (1) at (1.5, 1) {1};
                \node<3->[node_visited] (1) at (1.5, 1) {1};

                \node<1-4>[node_base, opacity=0.3] (3) at (1.5, -1) {3};
                \node<5>[node_current] (3) at (1.5, -1) {3};
                \node<6->[node_visited] (3) at (1.5, -1) {3};

                \node<1-2>[node_base, opacity=0.3] (2) at (3, 1) {2};
                \node<3>[node_current] (2) at (3, 1) {2};
                \node<4->[node_visited] (2) at (3, 1) {2};

                \node<1-3>[node_base, opacity=0.3] (4) at (3, -1) {4};
                \node<4>[node_current] (4) at (3, -1) {4};
                \node<5->[node_visited] (4) at (3, -1) {4};

                % Vẽ cạnh với hiệu ứng highlight
                \draw<2>[orange, ultra thick] (0)--(1); % Cạnh đang duyệt
                \draw<1,3->[gray!50] (0)--(1);          % Cạnh đã/chưa duyệt
                
                \draw<5>[orange, ultra thick] (0)--(3);
                \draw<1-4,6->[gray!50] (0)--(3);

                \draw<3>[orange, ultra thick] (1)--(2);
                \draw<1-2,4->[gray!50] (1)--(2);

                \draw<4>[orange, ultra thick] (2)--(4);
                \draw<1-3,5->[gray!50] (2)--(4);
            \end{tikzpicture}
        \end{column}
        \begin{column}{0.45\textwidth}
            % Giữ nguyên phần Stack của bạn hoặc thêm \pause để khớp nhịp
            \textbf{Trạng thái Stack:}\\
            \begin{tikzpicture}
                \draw[thick] (0,2.5) -- (0,0) -- (1.5,0) -- (1.5,2.5);
                \node<1-4>[stack_box] at (0.75, 0.35) {3};
                \node<1-1>[stack_box] at (0.75, 1.05) {1};
                \node<2-2>[stack_box] at (0.75, 1.05) {2};
                \node<3-3>[stack_box] at (0.75, 1.05) {4};
            \end{tikzpicture}
            \vspace{0.2cm} \small
            \only<1>{B1: Thăm 0, đẩy láng giềng \{3, 1\} vào Stack.}
            \only<2>{B2: Lấy 1 ra thăm. Thêm láng giềng 2 vào Stack.}
            \only<3>{B3: Lấy 2 ra thăm. Thêm láng giềng 4 vào Stack.}
            \only<4>{B4: Lấy 4 ra thăm. 4 không có láng giềng mới.}
            \only<5>{B5: Lấy 3 ra thăm. Ngăn xếp rỗng.}
        \end{column}
    \end{columns}
\end{frame}

% Slide 7: DFS Code
\begin{frame}[fragile]{Mã nguồn DFS (C++)}
\begin{verbatim}
#include <bits/stdc++.h>
using namespace std;

vector<int> adj;
bool visited;

void dfs(int u) {
    cout << u << " "; // Tham dinh u
    visited[u] = true;
    for (int v : adj[u]) {
        if (!visited[v]) {
            dfs(v); // De quy di sau vao nhanh v
        }
    }
}
\end{verbatim}
\end{frame}

\section{Thuật toán BFS}
% Slide 8: BFS Theory
\begin{frame}{4. BFS: Tìm kiếm theo chiều rộng}
    \begin{itemize}
        \item \textbf{Tư tưởng:} Duyệt theo từng "lớp" vết dầu loang. Thăm hết tất cả các đỉnh kề trực tiếp rồi mới tiến ra xa hơn.
        \item \textbf{Đặc điểm:} Luôn tìm được đường đi ngắn nhất trong đồ thị không trọng số.
        \item \textbf{Độ phức tạp:} Thời gian $O(V + E)$, Không gian $O(V)$.
    \end{itemize}
\end{frame}

% Slide 9: BFS Animation (Đã nâng cấp hiệu ứng)
\begin{frame}{Mô phỏng từng bước BFS}
    \begin{columns}
        \begin{column}{0.5\textwidth}
            \centering
            \begin{tikzpicture}[scale=1.2, every node/.style={node_base}]
                \node<1>[node_current] (0) at (0,0) {0};
                \node<2->[node_visited] (0) at (0,0) {0};

                \node<1>[node_base, opacity=0.3] (1) at (1.5, 1) {1};
                \node<2>[node_current] (1) at (1.5, 1) {1};
                \node<3->[node_visited] (1) at (1.5, 1) {1};

                \node<1-2>[node_base, opacity=0.3] (3) at (1.5, -1) {3};
                \node<3>[node_current] (3) at (1.5, -1) {3};
                \node<4->[node_visited] (3) at (1.5, -1) {3};

                \node<1-3>[node_base, opacity=0.3] (2) at (3, 1) {2};
                \node<4>[node_current] (2) at (3, 1) {2};
                \node<5->[node_visited] (2) at (3, 1) {2};

                \node<1-4>[node_base, opacity=0.3] (4) at (3, -1) {4};
                \node<5>[node_current] (4) at (3, -1) {4};
                \node<6->[node_visited] (4) at (3, -1) {4};

                % Hiệu ứng loang: Highlight tất cả cạnh kề khi thăm 0
                \draw<1>[blue, dashed, thick] (0)--(1) (0)--(3); 
                \draw<2->[gray!50] (0)--(1) (0)--(3);
                
                \draw<2>[blue, dashed, thick] (1)--(2);
                \draw<1,3->[gray!50] (1)--(2);

                \draw<4>[blue, dashed, thick] (2)--(4);
                \draw<1-3,5->[gray!50] (2)--(4);
            \end{tikzpicture}
        \end{column}
        \begin{column}{0.45\textwidth}
            % Phần Queue giữ nguyên logic của bạn
            \textbf{Hàng đợi Queue:}\\
            \begin{tikzpicture}
                \draw[thick] (0,0.7) -- (3,0.7) (0,0) -- (3,0);
                \node<1>[stack_box, minimum width=0.7cm] at (0.5, 0.35) {1};
                \node<1>[stack_box, minimum width=0.7cm] at (1.3, 0.35) {3};
                \node<2>[stack_box, minimum width=0.7cm] at (0.5, 0.35) {3};
                \node<2>[stack_box, minimum width=0.7cm] at (1.3, 0.35) {2};
                \node<3>[stack_box, minimum width=0.7cm] at (0.5, 0.35) {2};
                \node<4>[stack_box, minimum width=0.7cm] at (0.5, 0.35) {4};
            \end{tikzpicture}
            \vspace{0.2cm} \small
            \only<1>{B1: Thăm 0. Đưa láng giềng \{1, 3\} vào Queue.}
            \only<2>{B2: Lấy 1 ra thăm. Thêm láng giềng 2 vào Queue.}
            \only<3>{B3: Lấy 3 ra thăm. Đã duyệt 0.}
            \only<4>{B4: Lấy 2 ra thăm. Thêm láng giềng 4 vào Queue.}
            \only<5>{B5: Lấy 4 ra thăm. Hoàn tất.}
        \end{column}
    \end{columns}
\end{frame}

% Slide 10: BFS Code
\begin{frame}[fragile]{Mã nguồn BFS (C++)}
\begin{verbatim}
#include <bits/stdc++.h>
using namespace std;

void bfs(int start_node) {
    queue<int> q;
    q.push(start_node);
    visited[start_node] = true;

    while (!q.empty()) {
        int u = q.front(); q.pop();
        cout << u << " ";
        for (int v : adj[u]) {
            if (!visited[v]) {
                visited[v] = true;
                q.push(v); // Them cac dinh kề vao cuoi hang doi
            }
        }
    }
}
\end{verbatim}
\end{frame}

\section{Tổng kết}
% Slide 11: So sánh
\begin{frame}{So sánh DFS và BFS}
    \centering
    \begin{tabular}{|l|p{3.5cm}|p{3.5cm}|}
        \hline
        \textbf{Tiêu chí} & \textbf{DFS} & \textbf{BFS} \\ \hline
        \textbf{Cấu trúc} & Stack (Ngăn xếp) & Queue (Hàng đợi) \\ \hline
        \textbf{Đường đi} & Đi sâu (Mê cung) & Loang rộng (GPS) \\ \hline
        \textbf{Bộ nhớ} & Thấp hơn & Cao hơn \\ \hline
        \textbf{Ứng dụng} & Phát hiện chu trình, Topo & Tìm đường ngắn nhất \\ \hline
    \end{tabular}
\end{frame}
